\section{Conclusions}

    The aim of this thesis was to solve the proximal optimal consensus problem for a class of nonlinear agents with non-convex cost functions and constraints. In order to achieve consensus, a regulation form of Prox-GPDA, \cite{pmlr-v70-hong17a}, was designed with the addition of the NOn-Convex Gradient PROXimal (NOCGPROX) variables. These variables use sampled output measurements of the agents in a continuous manner to transform the optimal consensus problem into a regulation problem. In addition, robust distributed control laws were designed in order to achieve optimal consensus between the agents. The control laws were designed with a sliding mode backstepping approach with the use of Nussbaum gains. The crucial element for the applicability of this algorithm was the guarantee of the quadratic integration of the suggested variables. This approach enabled the algorithm to achieve sublinear convergence with a small number of iterations, while all signals of the closed-loop system proved to be bounded, attaining asymptotic stability. As a result, for the first time was proved that during the information exchange between the agents, needed only a finite amount of energy while transmitting the output sampled measurements. Therefore, the UAVs achieve optimal consensus with the proposed distributed control laws and algorithm.

\newpage
\section{Future Work}

    The proposed methodology that was researched in this thesis assumes that the output measurements are known. In certain application, sensors can not measure the output of a system. Thus, it is difficult to apply the above proposition. A sufficient solution to this problem is to seek Nash or Stackelberg equillibriums. This method does not require the knowledge of the output, hence it is estimated from the use of dynamic games. Another extension to this thesis is the use of switching topology. In that case, new robust distributed laws need to be applied to ensure that the switching of the topology does not have an impact in the communications between the agents. Therefore, consensus will be achieved. Another extension to the problem is the use of Control Barrier Functions to ensure a safe environment for the agents in the presence of obstacles. The inclusion of these constraints will create reliable planning for the agents to achieve optimal consensus regarding the position of these obstacles. Additionally, another area that this methodology can be tested is in networks with malicious agents. Due to the fact that the information exchange between the agents is occurring in lower band frequencies with a certain delay factor, the network can be expanded with these agents to increase the security of the network. As a result, new distributed control laws will be developed to guarantee the secure exchange of information between the agents. 